\section{\texorpdfstring{\code{send()}}{send()}}
\label{sec:send}

\begin{cppcode}
bool Spi::send(const Frame& frame)
{
    if (!isValid(frame))                              { return false; }
    if ((readReg(RegStatus) & StatusTxReady) == 0U)   { return false; }

    writeReg(RegTxId,      frame.id);
    writeReg(RegTxDlc,     frame.dlc);
    writeReg(RegTxDataLo,  packDataLo(frame));
    writeReg(RegTxDataHi,  packDataHi(frame));
    writeReg(RegTxSend,    Trigger);

    return true;
}
\end{cppcode}

\subsection{Ordningen är inte godtycklig}

\code{TX_SEND} \textbf{måste} skrivas sist. Den skrivningen betyder inte "spara ett värde" utan
"börja sända nu", och kontrollern läser då \code{TX_ID}, \code{TX_DLC} och de två dataregistren.
Skriv triggern först och du skickar förra framens innehåll --- vilket vid bring-upen ser ut som en
nod som ligger exakt en frame efter.

Statuskontrollen måste komma \textbf{före} skrivningarna. Skulle du skriva \code{TX_ID} medan en
sändning pågår skriver du över data kontrollern håller på att använda.

\subsection{Två \texorpdfstring{\code{false}}{false}, två betydelser}

Interfacet har bara en \code{bool} att uttrycka båda med:

\begin{narrowtable}{@{}ll@{}}
\thead{Orsak} & \thead{Vad anroparen bör göra} \\ \midrule
Ogiltig frame & Rätta framen. Väntan hjälper aldrig. \\
Hårdvaran upptagen & Vänta och försök igen. \\
\end{narrowtable}

Det är en medveten förenkling, men värd att vara medveten om. Två saker följer:

\textbf{Ordningen spelar roll.} Valideringen först, så att en ogiltig frame avvisas \emph{utan} att
en enda transaktion går ut på bussen. Den utdelade testsviten kontrollerar just det.

\textbf{Det blir ett problem exakt en söm högre upp.} Den här loopen är en oändlig loop för en
frame med \code{dlc = 9}:

\begin{cppcode}
while (!driver.send(frame))   // hänger för alltid om framen är ogiltig
{
}
\end{cppcode}

Rätt hantering är att validera innan man går in i loopen, eller att räkna försök och ge upp.

\section{Packningen}
\label{sec:packing}

Databyte 0 ligger i de \textbf{mest signifikanta} bitarna, och \code{TX_DATA_LO} bär de
\textbf{första} fyra bytesen.

\begin{cppcode}
/**
 * @brief Pack four payload bytes into one register word, byte 0 in the most significant bits.
 */
[[nodiscard]] std::uint32_t pack(const std::uint8_t* const data,
                                 const std::uint8_t offset) noexcept
{
    return (static_cast<std::uint32_t>(data[offset + 0U]) << 24U)
         | (static_cast<std::uint32_t>(data[offset + 1U]) << 16U)
         | (static_cast<std::uint32_t>(data[offset + 2U]) << 8U)
         | (static_cast<std::uint32_t>(data[offset + 3U]));
}
\end{cppcode}

Samma cast-regel som i \kapref{ch:byte}: \textbf{casta före skiftningen}.

De två vanligaste packningsfelen:

\begin{booktable}{@{}lL@{}}
\thead{Fel} & \thead{Symptom} \\ \midrule
\code{TX_DATA_LO} och \code{TX_DATA_HI} förväxlade & Framen kommer fram med byte 4--7 först.
  Syns direkt i en bussanalysator, och inte alls i simulering om testet också har fel. \\
Skiftningarna i omvänd ordning & \code{AA BB CC DD} kommer fram som \code{DD CC BB AA}. \\
\end{booktable}

Båda ger en frame som \emph{kommer fram} --- bara med fel innehåll. Det är den obehagligaste
sortens fel, för allt \emph{ser} ut att fungera.

\section{\texorpdfstring{\code{receive()}}{receive()}}
\label{sec:receive}

\begin{cppcode}
bool Spi::receive(Frame& frame)
{
    if ((readReg(RegStatus) & StatusRxValid) == 0U) { return false; }

    frame.id  = static_cast<std::uint16_t>(readReg(RegRxId) & 0x7FFU);
    frame.dlc = static_cast<std::uint8_t>(readReg(RegRxDlc) & 0xFU);

    unpack(readReg(RegRxDataLo), frame.data, 0U);
    unpack(readReg(RegRxDataHi), frame.data, 4U);

    for (std::uint8_t i{frame.dlc}; i < Frame::MaxDataLen; ++i)
    {
        frame.data[i] = 0U;
    }

    writeReg(RegRxAck, Trigger);
    return true;
}
\end{cppcode}

\subsection{Maskeringen mot \texorpdfstring{\code{dlc}}{dlc}}

Loopen ser överflödig ut. Registerkartan säger ju att kontrollern nollställer sin dataackumulator
vid varje framestart, så de bytesen \emph{är} redan noll.

Gör det ändå, av tre skäl:

\begin{enumerate}
  \item \textbf{Nollorna är kontrollerns garanti, inte framens innehåll.} En drivrutin som
    kopierar åtta bytes rakt av rapporterar upp till sex bytes som aldrig sändes. Att de råkar
    vara noll gör dem inte till data.
  \item \textbf{Garantin kan försvinna.} Den kommer från en implementationsdetalj i hårdvarans
    VHDL. Ändras den, eller lånar du ett kort från en annan grupp, är din drivrutin fortfarande
    rätt.
  \item \textbf{\code{dlc} är kontraktet.} \code{frame.dlc = 2} betyder "två bytes gäller". Kod
    som lämnar ifrån sig ett objekt vars fält säger emot varandra gör nästa läsare osäker.
\end{enumerate}

\subsection{\texorpdfstring{\code{RX_ACK}}{RX\_ACK} sist}

Kvittensen frigör \code{RX_*}-registren och nollställer \code{STATUS} bit 1. Skriv den
\textbf{efter} att du läst allt --- annars kan nästa frame skriva över registren mitt i din
läsning.

Glöm den helt, och bit 1 står kvar för alltid: varje \code{receive()} returnerar \code{true} med
samma frame, i en oändlig loop. Det är ett symptom värt att känna igen.

\section{\texorpdfstring{\code{hasError()}}{hasError()} och \texorpdfstring{\code{clearError()}}{clearError()}}
\label{sec:errors}

\begin{cppcode}
bool Spi::hasError() const noexcept
{
    return (readReg(RegStatus) & StatusError) != 0U;
}

void Spi::clearError() noexcept
{
    writeReg(RegErrorFlags, 0U);
}
\end{cppcode}

Två rader, och en av dem är en fälla.

\textbf{\code{clearError()} skriver \code{0x0}, inte \code{0x1}.} \code{ERROR_FLAGS} är en
lås-och-nollställ-på-noll-konstruktion: bara en skrivning där hela ordet är noll nollställer
låsningen. \code{0x1} gör ingenting, och \code{0x000000FF} gör ingenting trots att bit 0 är satt.

Det är motsatt regel mot \code{TX_SEND} och \code{RX_ACK}, som kräver att bit 0 \emph{är} satt.
Att de tre registren har olika regler är inte elegant, men det står i kartan, och det är kartan
som gäller.

\textbf{Symptomet om du skriver \code{0x1}:} felbiten går aldrig ned, \code{hasError()} returnerar
\code{true} för alltid efter första felet, och all felhantering ovanför slutar fungera.

\section{Att vänta på en sändning}
\label{sec:waiting}

\code{send()} \textbf{väntar inte}. Den returnerar så fort triggern är skriven, och det är rätt: en
drivrutinsmetod som blockerar i upp till 130~µs på en 4~MHz-MCU äter hela systemets tidsbudget.

Hur väntar anroparen då? Interfacet har varken \code{txReady()} eller \code{waitForSend()} ---
avsiktligt, eftersom \code{Stub} inte kunnat implementera dem meningsfullt. Det anroparen har är
att \code{send()} returnerar \code{false} medan sändaren är upptagen. \textbf{Väntan är alltså ett
nytt försök:}

\begin{cppcode}
driver.clearError();

if (!isValid(frame)) { return false; }   // never going to be accepted, however long we wait

std::uint16_t attempts{};
while (!driver.send(frame))
{
    if (++attempts >= MaxAttempts) { return false; }
}

// The port was free when the frame was accepted, so the previous transmission had finished.
if (driver.hasError())
{
    driver.clearError();
    // Lost arbitration, stuff error, bad CRC or a remote frame - the latch does not say which.
}
\end{cppcode}

\begin{limitation}
\textbf{En förlorad arbitrering ser ut som en lyckad sändning.} \code{STATUS} bit 0 kommer tillbaka
i båda fallen, och betyder \emph{porten är fri}, inte \emph{framen kom fram}. Vill du veta
utfallet måste du läsa \code{ERROR_FLAGS} \emph{efter} att biten kommit tillbaka --- och även då
får du veta \emph{att} något gick fel, aldrig \emph{vad}.
\end{limitation}

\section{Checklista}
\label{sec:checklist}

\begin{itemize}
  \item[\checkbox] \code{send()} validerar \textbf{först}, och skickar ingen transaktion för en
    ogiltig frame.
  \item[\checkbox] \code{send()} kontrollerar \code{StatusTxReady} innan den skriver något.
  \item[\checkbox] \code{TX_SEND} skrivs \textbf{sist}.
  \item[\checkbox] \code{TX_DATA_LO} bär databyte 0--3, med byte 0 i bitarna 31--24.
  \item[\checkbox] Varje \code{static_cast<std::uint32_t>} står \textbf{före} skiftningen.
  \item[\checkbox] \code{receive()} kontrollerar \code{StatusRxValid} först.
  \item[\checkbox] \code{receive()} skriver \code{RX_ACK} \textbf{efter} alla läsningar.
  \item[\checkbox] \code{receive()} nollställer \code{frame.data} ovanför \code{frame.dlc}.
  \item[\checkbox] \code{clearError()} skriver \code{0x0}.
  \item[\checkbox] Ingen magisk siffra i hela filen.
\end{itemize}
